\documentclass[11pt]{article}

\usepackage[margin=2.5cm]{geometry}
\usepackage{amsmath, amssymb}
\usepackage{booktabs}
\usepackage{hyperref}
\usepackage{enumitem}
\usepackage{xcolor}

\newcommand{\madsen}[1]{\textcolor{gray}{[Madsen, \S#1]}}

\title{Understanding the Karno tank signals: a statistical route\\
\large following Madsen, \emph{Time Series Analysis}}
\author{}
\date{\today}

\begin{document}
\maketitle

\section{Context}

The grey-box model (\texttt{scripts/graybox.R}, \texttt{ctsmTMB}) treats the
storage tank as two lumped thermal masses, $T_{\text{top}}$ and
$T_{\text{bot}}$, coupled by a linear conductive term $k\,(T_{\text{top}} -
T_{\text{bot}})$, driven by the air-source heat pump's electrical power
$P_{\text{el}}$ (into the top node) and the district-heating draw
$\dot Q_{\text{dist}}$ (at the bottom node), with a loss term $UA\,(T -
T_{\text{amb}})$ on each node.

Fitting this model to $\sim$11 days of real logged data
(\texttt{data/karno-410708\_raw\_k0001.parquet}, 2026-05-13 to 2026-05-24,
5-minute sampling) produces a badly behaved fit: several parameters
($C_{\text{top}}$, $C_{\text{bot}}$, $k$, $a$) pin at whatever bound they are
given rather than converging to an interior optimum, and the measurement/process
noise parameters split asymmetrically between the two channels ($\sigma_{y_1}
\to$ its upper bound, $\sigma_{y_2} \to$ its lower bound, or vice versa
depending on how the bounds are set) --- the classic sign of the optimizer
using one channel's noise term to paper over a structural mismatch instead of
a real fit.

This document is not a fix. It is a statistical route --- grounded in Madsen's
\emph{Time Series Analysis} --- for actually characterizing the input/output
dynamics of the two sensors \emph{before} committing to any particular ODE
structure again.

\subsection{Physical layout}

\begin{center}
\begin{tabular}{c}
top of tank \\
$\vert$ 30 cm \\
$T_{\text{top}}$ sensor \\
$\vert$ 70 cm \\
$T_{\text{bot}}$ sensor \\
$\vert$ 30 cm \\
bottom of tank
\end{tabular}
\end{center}

Heat is added at the $T_{\text{top}}$ level; the district-heating draw
(``cool'') is taken out at the $T_{\text{bot}}$ level. Despite this
asymmetric but plausible geometry, the observed behaviour is not what a
simple stratified/conductive model predicts:

\begin{itemize}
\item $T_{\text{top}}$: flat, then declines with no heat-pump input; rises
  when the heat pump runs. This is expected -- it is the injection point.
\item $T_{\text{bot}}$: rises \emph{directly} when the heat pump runs, with
  no visible delay consistent with a thermocline having to migrate 70\,cm
  down past the sensor.
\end{itemize}

\section{The route}

The steps below are ordered the way they should actually be executed: each
step's output is required input for the next. Section numbers refer to
Madsen, \emph{Time Series Analysis}.

\subsection*{Step 0 --- univariate structure} \madsen{5.5, 6.2.1}

Before touching cross-relationships, look at each series alone:
autocorrelation function (ACF) and partial ACF (PACF) of $T_{\text{top}}$,
$T_{\text{bot}}$, $P_{\text{el}}$, $\dot Q_{\text{dist}}$ individually
(\S6.2.1, estimation of the ACF; \S5.5 for what AR/MA signatures look like in
the ACF/PACF). At 5-minute sampling all four series will be highly
persistent (ACF decaying slowly, near unit-root behaviour). This matters
directly for Step 2: it is exactly what makes a \emph{raw} cross-correlation
between two persistent series unreliable.

\subsection*{Step 1 --- raw cross-correlation function (CCF)} \madsen{6.2.2}

The cross-covariance and cross-correlation functions,
\[
  \hat\gamma_{xy}(\ell) = \frac{1}{n}\sum_{t} (x_t - \bar x)(y_{t+\ell} -
  \bar y), \qquad
  \hat\rho_{xy}(\ell) = \frac{\hat\gamma_{xy}(\ell)}{\sqrt{\hat\gamma_{xx}(0)\,
  \hat\gamma_{yy}(0)}},
\]
computed for a range of lags $\ell$, is the formal version of ``lag
correlation''. Do this first as a baseline, but treat the result with
suspicion: two autocorrelated series will show spuriously significant CCF
values at many lags purely because of their own internal persistence, not
because one drives the other. This is why Step 2 exists.

\subsection*{Step 2 --- pre-whitening} \madsen{9.6.1}

The standard (Box--Jenkins) fix for the Step 1 problem:
\begin{enumerate}[label=(\roman*)]
\item Fit a model (AR($p$) is enough in practice) to the input series $x_t$
  such that the residuals $a_t = x_t - \sum_i \phi_i x_{t-i}$ are
  approximately white noise.
\item Apply the \emph{same} filter to the output series: $b_t = y_t -
  \sum_i \phi_i y_{t-i}$.
\item Compute the CCF between $a_t$ and $b_t$.
\end{enumerate}
Because $a_t$ is white, the resulting cross-correlation is a clean,
interpretable estimate of the dynamic relationship between $x$ and $y$ --
its shape approximates the system's impulse response (Step 3), and its
confidence bands are the simple $\pm 1.96/\sqrt{n}$ used for white-noise
input.

\emph{This has already been run} (\texttt{scripts/diagnostics.R}) for
$P_{\text{el}} \to T_{\text{top}}$ and $P_{\text{el}} \to T_{\text{bot}}$.
Results, with the convention that a positive lag means the input leads the
output:

\begin{center}
\begin{tabular}{lccc}
\toprule
Pair & AR order & peak lag & peak $r$ \\
\midrule
$P_{\text{el}} \to T_{\text{bot}}$ & 20 & $+1$ sample ($+5$ min) & $+0.296$ \\
$P_{\text{el}} \to T_{\text{top}}$ & 20 & $0$ samples ($0$ min) & $-0.393$ \\
\bottomrule
\end{tabular}
\end{center}

$P_{\text{el}} \to T_{\text{bot}}$ is exactly what a direct heat input
should look like: small positive lag, positive sign. $P_{\text{el}} \to
T_{\text{top}}$ is not: the strongest relationship is \emph{instantaneous
and negative}, at no lag where a delayed positive response would be
expected from direct heat injection at that sensor's level.

\subsection*{Step 3 --- reading the CCF as an impulse response} \madsen{8.3--8.4}

Once pre-whitened, the shape of the cross-correlation function \emph{is} a
(scaled) estimate of the system's impulse-response weights $v_0, v_1, v_2,
\dots$ relating input to output (\S8.3.1, transfer function models;
\S8.4, identification of transfer function models from the CCF). In
practice: a single sharp spike at some lag $\ell_0$ suggests a simple gain
plus dead time of $\ell_0$ steps; a spike that decays geometrically over
several lags suggests a first-order lag (single time constant); an
oscillating or sign-changing pattern (which is what both the raw data and
the model residuals show here, see Step 5) suggests either an
under-differenced/near-unit-root series contaminating the estimate, or --
more likely given the context -- a genuinely mis-specified relationship
that no simple transfer function order will resolve, because the assumed
causal direction is wrong (see \S\ref{sec:caveat}).

\subsection*{Step 4 --- multiple-input extension} \madsen{8.5}

There are four real candidate inputs ($P_{\text{el}}$, $\delta$ = GSHP
on/off, $\dot Q_{\text{dist}}$, $T_{\text{amb}}$). \S8.5 covers identifying
and estimating transfer functions with several regressors simultaneously
(moment relations \S8.5.1, spectral relations \S8.5.2, identification
\S8.5.3) rather than one input at a time, which matters once individual
pairwise CCFs are understood and need to be combined without
double-counting shared variance (e.g.\ $P_{\text{el}}$ and $\delta$ are
correlated with each other, since both indicate heat-pump activity).

\subsection*{Step 5 --- residual analysis of the currently fitted model} \madsen{6.6}

For the model that is already fitted (\texttt{fit} object from
\texttt{model\$estimate()} in \texttt{ctsmTMB}), the one-step-ahead
standardized innovations (\texttt{fit\$residuals\$normalized}) should behave
like white noise \emph{if} the state equations correctly capture the
system's dynamics. Two checks (\S6.6.2, residual analysis):
\begin{enumerate}[label=(\alph*)]
\item ACF of each residual series alone -- significant autocorrelation means
  the state's own dynamics (not just the inputs) are under-specified.
\item Cross-correlation of each (pre-whitened, Step 2) input against the
  residual series -- significant correlation at some lag means that input's
  effect is not fully captured by the fitted model at that lag.
\end{enumerate}

\emph{Already run.} Result: both residual series show strong, alternating-sign,
slowly-decaying autocorrelation out past lag 40 (200 minutes) -- the current
fit is a poor one-step predictor for both states, not just $T_{\text{top}}$.
However, cross-correlating the residuals against $P_{\text{el}}$ shows
\emph{no} significant lag for either channel. That combination is
informative: the problem is not a missing lagged copy of $P_{\text{el}}$ in
the state equation (which residual cross-correlation would have caught) --
it is something the current noise-parameter degeneracy (Step-2 finding) is
already masking. Re-run this check once the noise-parameter pathology is
resolved (i.e.\ interior, non-bound-pinned $\sigma$ estimates), since a
degenerate Kalman gain distorts residual whiteness independently of any
missing input dynamics.

\subsection*{Step 6 --- frequency domain cross-check} \madsen{7.3}

As a complement to the time-domain CCF: the cross-spectrum
(\S7.3) decomposes the input/output relationship by frequency. The
\emph{phase spectrum} (\S7.3.2) gives lag as a function of frequency
(useful if the delay itself is frequency-dependent, e.g.\ short bursts of
heat-pump operation propagate differently than long ones), and the
\emph{coherence spectrum} gives, per frequency, how much of the output's
variance the input actually explains -- useful here because $\delta$ (GSHP
on/off) is a sparse, near-impulsive signal (duty cycle $\approx 0.5\%$,
5 transitions in 11 days) while $P_{\text{el}}$ is closer to continuous; the
two may show up very differently in a coherence plot.

\subsection*{Step 7 --- the worked example in the book} \madsen{1.1.3}

Section 1.1.3, ``Heat dynamics of a building'', is the book's own example of
grey-box RC-network modeling of a thermal system -- structurally the same
model class used here. Worth reading directly for how the book frames
identifiability and model-order selection for this kind of problem.

\section{A caveat outside the book's scope}
\label{sec:caveat}

Steps 1--6 (and the whole of Chapter 8) implicitly assume the input is
\emph{exogenous} -- something that happens independently of the state being
modeled. The Step-2 result above (instantaneous, negative
$P_{\text{el}} \to T_{\text{top}}$ correlation, versus the physically sane
delayed positive $P_{\text{el}} \to T_{\text{bot}}$ correlation) is not
consistent with an exogenous input at all: it is the signature of
\emph{closed-loop} behaviour, where $P_{\text{el}}$ is itself a function of
$T_{\text{top}}$ via the heat pump's thermostat (power ramps up when
$T_{\text{top}}$ is low, throttles down as it recovers). This was confirmed
directly: the heat pump's on/off command \texttt{z\_HP,air} correlates
$0.95$ with $P_{\text{el}}$ (expected) but \emph{also} $-0.48$ with
$T_{\text{top}}$ -- the command itself is triggered by low $T_{\text{top}}$,
and the setpoint \texttt{T\_set,HP,air} is constant ($=0$) in this window, so
there is no clean alternative exogenous proxy available here.

Madsen's book, being a classical (largely open-loop) time series text, does
not cover closed-loop identification (see e.g.\ Ljung, \emph{System
Identification: Theory for the User}, for direct/indirect closed-loop
methods). What its tools \emph{did} do is surface the problem cleanly --
the wrong-signed, zero-lag CCF peak is exactly what closed-loop feedback
looks like under this diagnostic. Any structural fix (modeling the
thermostat explicitly, restricting to genuinely open-loop stretches of the
log, or treating $T_{\text{top}}$ as a measured boundary condition rather
than a state to be solved for) is a separate design decision, deliberately
left open here.

\section{Practical checklist}

\begin{enumerate}
\item Run Step 0 (ACF/PACF) on all four series to confirm the persistence
  that motivates pre-whitening.
\item Extend \texttt{scripts/diagnostics.R}'s pre-whitened CCF (already
  covers $P_{\text{el}}$, $\dot Q_{\text{dist}}$) to $\delta$ and
  $T_{\text{amb}}$ as well.
\item Read off Step 3 (impulse-response shape) for whichever pairs come back
  with a clean, physically sane sign and lag -- $T_{\text{bot}}$'s channel
  looks like the promising one.
\item Re-run Step 5 (residual analysis) once the model's noise parameters
  are no longer bound-pinned, to separate genuine missing dynamics from
  Kalman-filter noise-mismatch artifacts.
\item Run Step 6 (coherence/phase spectrum) as a second opinion on Step 2,
  particularly for $\delta$ given how sparse it is.
\item Treat $P_{\text{el}}$ (and $\delta$/\texttt{z\_HP,air}) as
  provisionally endogenous until/unless an open-loop stretch of data or an
  explicit controller model says otherwise.
\end{enumerate}

\end{document}
